% sec_discussion.tex -- Section 9: the 7x7x7 board, remarks and open problems (editor)
\section{The \texorpdfstring{$7\times7\times7$}{7x7x7} board, remarks and open
problems}\label{sec:discussion}

\subsection{The \texorpdfstring{$7\times7\times7$}{7x7x7} board}\label{ssec:seven}

\Cref{thm:main} implies the corresponding statement for the board of side~$7$. This statement was first
proved by Cain, Raymond and K\"allersj\"o~\cite{CainRaymondKallersjo2025}, by a direct computation on
the $7\times7\times7$ board in which pairings are computed for about $3.5\cdot 10^7$ positions. The
proof given here derives it from \cref{thm:main} by an argument that can be checked by hand. It rests
on the decomposition of a board of odd side by its mid-planes, which Golomb and Hales use in the proof
of~\cite[Theorem~3]{GolombHales2002}.

The $7\times7\times7$ board is $\Qseven=\{0,\dots,6\}^3$. Its lines are defined as in \cref{ssec:lines}:
a \emph{parametrization} is a triple $f=(f_1,f_2,f_3)$ of maps $\{0,\dots,6\}\to\{0,\dots,6\}$, each of
which is constant, the identity or the reversal $\rho_7\colon i\mapsto 6-i$, and which are not all
constant; a component that is not constant is called \emph{active}; and the line of $f$ is
$\ell(f)=\{f(t):0\le t\le 6\}$. It consists of seven distinct cells, because an active component is
injective, and $\ell(f\circ\rho_7)=\ell(f)$. There are $193$ lines~\cite[Theorem~1]{GolombHales2002}.
Let
\[
  M=\{c\in\Qseven : c_i=3 \text{ for some } i\},\qquad A=\Qseven\setminus M=\{0,1,2,4,5,6\}^3 .
\]
Thus $M$ is the union of the three mid-planes, $|A|=216$ and $|M|=343-216=127$. Let $\varphi$ be the
increasing bijection from $\{0,1,2,4,5,6\}$ onto $\{0,\dots,5\}$, that is, $\varphi(i)=i$ for $i\le 2$
and $\varphi(i)=i-1$ for $i\ge 4$, and let
\[
  \Phi\colon A\to\Q,\qquad \Phi(c)=(\varphi(c_1),\varphi(c_2),\varphi(c_3)).
\]
Note that $\varphi(6-i)=5-\varphi(i)$ for all $i\ne 3$.

\begin{lemma}[Mid-plane decomposition]\label{lem:seven-decomp}
Let $f$ be a parametrization of a line of $\Qseven$.
\begin{enumerate}[label=\textup{(\alph*)}]
\item If some component of $f$ is the constant map with value $3$, then $\ell(f)\subseteq M$.
\item Otherwise $\ell(f)\cap M=\{f(3)\}$, and $\Phi(\ell(f)\cap A)$ is a line of the $6\times6\times6$
      board~$\Q$.
\item Every line of $\Q$ is equal to $\Phi(\ell\cap A)$ for exactly one line $\ell$ of $\Qseven$.
\item The lines of $\Qseven$ that are contained in $M$ are the $45$ lines listed in
      \cref{tab:seven-pairing}.
\end{enumerate}
\end{lemma}

\begin{proof}
(a) is clear. (b) Suppose that no component of $f$ is the constant $3$. An active component takes the
value $3$ for $t=3$ and for no other $t$, and at least one component is active. Hence $f(t)\in M$ if and
only if $t=3$. For $t\ne 3$ put $t'=\varphi(t)$, which runs through $\{0,\dots,5\}$. The $i$-th
coordinate of $\Phi(f(t))$ is $\varphi(v)$ if $f_i$ is the constant $v$, it is $\varphi(t)=t'$ if $f_i$ is
the identity, and it is $\varphi(6-t)=5-t'$ if $f_i=\rho_7$. So $\Phi(\ell(f)\cap A)=\ell(f')$, where $f'$
is the parametrization, in the sense of \cref{ssec:lines}, that arises from $f$ by replacing every
constant $v$ by $\varphi(v)$, the identity by the identity, and $\rho_7$ by $\rho$.

(c) The map $f\mapsto f'$ is a bijection from the parametrizations of $\Qseven$ without a constant
component $3$ onto the parametrizations of $\Q$, and $(f\circ\rho_7)'=f'\circ\rho$. So every line of
$\Q$ arises as in (b). If $\ell(f')=\ell(g')$, then $g'\in\{f',f'\circ\rho\}$ by
\cref{lem:geometry}\,(a); hence $g\in\{f,f\circ\rho_7\}$ and $\ell(g)=\ell(f)$. A line $\ell(g)$ as in
(a) does not occur here, because $\ell(g)\cap A=\emptyset$.

(d) By (a) and (b) a line is contained in $M$ if and only if it has a parametrization with a constant
component $3$. Since $\ell(f)=\ell(f\circ\rho_7)$, we may assume that the first active component of $f$
is the identity. If exactly one component is active, the other two are constants $v,w$ with $v=3$ or
$w=3$, which gives $13$ triples for each of the three positions of the active component. If two
components are active, the third is the constant $3$ and the second active component is the identity
or $\rho_7$, which gives two triples for each of the three positions of the constant. These are the
$39+6=45$ triples in the first, third and fifth column of \cref{tab:seven-pairing}, where a triple is
written as its value at $t$.
\end{proof}

In particular $148$ of the $193$ lines meet $A$ in a line of the $6\times6\times6$ array, and
$193-148=45$ lines lie in $M$, in agreement with (d).

\begin{lemma}[A pairing for the mid-planes]\label{lem:seven-pairing}
\Cref{tab:seven-pairing} assigns to every line $\ell\subseteq M$ a set $\pi(\ell)$ of two cells of
$\ell$, and the $45$ sets $\pi(\ell)$ are pairwise disjoint.
\end{lemma}

\begin{proof}
This can be checked directly in the table; the following description makes the check short. Let
$\psi=\varphi^{-1}$, and read the indices $r,s\in\{0,\dots,5\}$ modulo $6$. In the mid-plane $z=3$ the
table assigns
\begin{itemize}
\item to the line $\{(\psi(r),t,3)\}$ the cells $(\psi(r),\psi(r+1),3)$ and $(\psi(r),\psi(r+2),3)$;
\item to the line $\{(t,\psi(s),3)\}$ the cells $(\psi(s+1),\psi(s),3)$ and $(\psi(s+2),\psi(s),3)$;
\item to the diagonal $\{(t,t,3)\}$ the cells $(0,0,3)$ and $(6,6,3)$, and to the diagonal
      $\{(t,6-t,3)\}$ the cells $(1,5,3)=(\psi(1),\psi(4),3)$ and $(5,1,3)=(\psi(4),\psi(1),3)$.
\end{itemize}
Each of these cells lies on the line to which it is assigned. For a cell $(\psi(r),\psi(s),3)$ consider
the residue of $s-r$ modulo $6$: it is $1$ or $2$ for the cells of the first kind, $5$ or $4$ for those
of the second kind, $0$ for the two cells of the first diagonal and $3$ for the two cells of the second
one. So cells assigned to lines of different kinds are different. The two cells of a line of the first
kind differ, and cells of different lines of the first kind have different values of $r$; likewise for
the second kind. The four cells assigned to the diagonals are distinct. Hence the $28$ cells assigned
in the plane $z=3$ are distinct. The assignments in the planes $y=3$ and $x=3$ are obtained from this
one by moving the coordinate $3$ to the second and to the first position, respectively. All cells
assigned so far have exactly one coordinate equal to $3$, so cells assigned in different mid-planes are
different. Finally the three lines $\{(t,3,3)\}$, $\{(3,t,3)\}$, $\{(3,3,t)\}$ receive the cells with
$t=0$ and $t=1$; these six cells are distinct and have two coordinates equal to $3$.
\end{proof}

\begin{corollary}[computer-assisted, via \cref{thm:main}]\label{cor:seven}
In the Maker--Breaker game on the $193$ lines of the $7\times7\times7$ board, \Breaker, moving second,
has a winning strategy: a strategy that prescribes an empty cell whenever \Breaker{} is to move and the
board is not full, and under which \Maker{} never owns all seven cells of a line. Consequently, ordinary
$7\times7\times7$ tic-tac-toe is a draw: each of the two players has a strategy with which this player
never loses.
\end{corollary}

\begin{proof}
By \cref{thm:main} there is a strategy $\sigma$ for \Breaker{} as the second player of the Maker--Breaker
game on $\Q$ that prescribes an empty cell whenever \Breaker{} is to move and under which \Maker{} never
owns all six cells of a line of $\Q$. We identify $A$ with $\Q$ by means of $\Phi$, and we use the
pairing $\pi$ of \cref{lem:seven-pairing}; the two cells of $\pi(\ell)$ are \emph{partners} of each
other.

\smallskip\noindent\emph{The strategy.} While playing on $\Qseven$, \Breaker{} keeps track of a
\emph{virtual play} of the game on $A$, in which the virtual \Breaker{} follows $\sigma$. Initially no
move has been made in the virtual play. Suppose that \Maker{} has just occupied a cell $x$ of $\Qseven$.
\begin{enumerate}[label=\textup{(\roman*)}]
\item If $x\in A$, the virtual \Maker{} occupies $x$, and the virtual \Breaker{} occupies the cell $w\in A$
      that $\sigma$ prescribes. If $w$ is empty on the real board, \Breaker{} occupies $w$; otherwise he
      makes a free move.
\item If $x\in M$, the virtual play is not continued. If $x\in\pi(\ell)$ for a line $\ell\subseteq M$ and
      the partner of $x$ is empty, \Breaker{} occupies the partner of $x$; otherwise he makes a free move.
\end{enumerate}
A \emph{free move} is an arbitrary empty cell of $\Qseven$. If there is none, the board is full and the
game is over.

\smallskip\noindent\emph{The invariant.} We claim that the following holds whenever \Maker{} is to move.
\begin{enumerate}[label=\textup{(J\arabic*)}]
\item The virtual play is a play of the game on $A$ in which \Breaker{} has followed $\sigma$ and
      \Maker{} is to move. The stones of the virtual \Maker{} are exactly the stones of \Maker{} in $A$,
      and every stone of the virtual \Breaker{} is a stone of \Breaker.
\item Every line $\ell\subseteq M$ contains a stone of \Breaker, or both cells of $\pi(\ell)$ are empty.
\end{enumerate}
Both statements hold at the beginning. Assume that they hold, and let \Maker{} occupy the cell $x$.

Let $x\in A$. By (J1) every cell of $A$ that is occupied in the virtual position is occupied on the
real board; so $x$ is empty in the virtual position, and the move of the virtual \Maker{} is legal. By
\cref{thm:main} the cell $w$ prescribed by $\sigma$ is empty in the virtual position. In particular $w$
is not a stone of \Maker: the stones of \Maker{} in $A$, including $x$, are the stones of the virtual
\Maker. Hence, if $w$ is not empty on the real board, it holds a stone of \Breaker. In either case $w$ is
a stone of \Breaker{} after the reply, and the reply adds no stone of \Maker. So (J1) holds again.

Let $x\in M$. Then the virtual play and the stones of \Maker{} in $A$ are unchanged, and the reply of
\Breaker{} adds at most a stone of \Breaker{} in $A$; so (J1) holds again.

For (J2) let $\ell\subseteq M$ and let $z$ be the reply of \Breaker{} (if the board is full after the
move $x$, there is nothing to show). If $\ell$ contained a stone of \Breaker, it still does. Otherwise
both cells of $\pi(\ell)$ were empty. If $x\in\pi(\ell)$, then $x\in M$ and the partner of $x$ is empty,
so $z$ is this partner, and $\ell$ now contains the stone $z$ of \Breaker. If $x\notin\pi(\ell)$, then
either $z\in\pi(\ell)\subseteq\ell$, and $\ell$ contains a stone of \Breaker, or both cells of
$\pi(\ell)$ are still empty.

\smallskip\noindent\emph{\Maker{} never owns a line.} The strategy prescribes an empty cell whenever
\Breaker{} is to move and the board is not full. Suppose that after some move $x$ of \Maker{} all seven
cells of a line $\ell$ of $\Qseven$ are stones of \Maker. If $\ell\subseteq M$, then $\ell$ contains no
stone of \Breaker, so by (J2), applied just before the move $x$, both cells of $\pi(\ell)$ were empty at
that moment; but only the cell $x$ of $\ell$ was empty then. If $\ell\not\subseteq M$, then by
\cref{lem:seven-decomp}\,(b) the six cells of $\ell\cap A$ form a line of $A=\Q$. By (J1), applied just
before the move $x$, and by rule (i) if $x\in A$, these six cells are stones of the virtual \Maker{} in
a play in which \Breaker{} follows $\sigma$. This contradicts the choice of $\sigma$.

\smallskip\noindent\emph{The ordinary game.} The ordinary game on $\Qseven$ is a game as in
\cref{prop:stealing}. If the second player follows the strategy just described, the first player never
owns a line; so the second player has a strategy with which he never loses, and the first player has no
winning strategy. The second player has no winning strategy by \cref{prop:stealing}. As in the proof of
\cref{cor:draw}, it follows from the determinacy of finite games~\cite{Zermelo1913} that the first
player also has a strategy with which he never loses.
\end{proof}

\begin{remark}\label{rem:seven}
(i) The free moves in the proof are the reason for the virtual play. For instance, the set $M$ has an
odd number of cells; when \Maker{} occupies its last empty cell and the board is not full, \Breaker{}
has to place a stone in $A$ that $\sigma$ has not asked for. The invariant (J1) expresses the well-known fact that an additional stone never
harms \Breaker: the virtual play ignores such a stone until $\sigma$ prescribes its cell, and \Maker{}
can never occupy it.

(ii) The proof uses \cref{thm:main} only through its statement, and \cref{lem:seven-decomp,lem:seven-pairing}
are verified by hand. The two lemmas were also confirmed by a short program
(\fn{scripts/check_7cube_midplanes.py} in the project repository), which enumerates the $193$ lines,
checks the statements (b)--(d) of \cref{lem:seven-decomp} and finds a pairing of the $45$ lines in $M$ by
bipartite matching. The pairing of \cref{tab:seven-pairing} consists of a pairing of the $14$ lines of
a $6\times6$ board, used in each of the three mid-planes, and of two cells on each of the three lines
through the center.

(iii) The outcomes for different $n$ are not known to be monotone: that a draw on the board $n^d$
implies a draw on the board $(n+1)^d$ is stated as a hypothesis by Golomb and
Hales~\cite[Section~8]{GolombHales2002}. The proof above establishes this implication for $d=3$ and
$n=6$, in the stronger form for \Breaker's win. It gives no information on the $5\times5\times5$
board.
\end{remark}

\subsection{Remarks}\label{ssec:remarks}

\subsubsection*{A Strong Draw that is not a pairing draw}
\Cref{thm:main} asserts more than a draw of the ordinary game: \Breaker{} wins the Maker--Breaker game as
the second player, which is a Strong Draw in the terminology of Beck~\cite{Beck2008}. On the other hand
the draw cannot be obtained by a pairing strategy, since no pairing exists on the empty board; by
\cref{prop:seven} a pairing strategy cannot even take over before \Breaker's seventh move. The
$6\times6\times6$ game thus belongs to the region in which Golomb and Hales placed it conjecturally:
the second player can force a draw, but not by a pairing strategy~\cite[Table~1]{GolombHales2002}. By
\cref{cor:seven} the same holds for the $7\times7\times7$ game.

\subsubsection*{The Erd\H{o}s--Selfridge criterion}
By the theorem of Erd\H{o}s and Selfridge~\cite{ErdosSelfridge1973}, in the form stated
in~\cite[Theorem~4.2]{Krivelevich2014}, \Breaker{} wins a Maker--Breaker game as the second player if the
winning sets $W$ satisfy $\sum_W 2^{-|W|}<1/2$. For the $6\times6\times6$ board the sum is
$148/64=2.3125$, so the criterion does not apply. Neither does its variant for winning sets of equal size
$k$ in which every cell lies in at most $\Delta$ of them, which requires that the number of winning
sets plus $\Delta$ is less than $2^k$ (quoted from~\cite[Theorem~6.1]{CainRaymondKallersjo2025}): here
$148+7=155\ge 64$. The sums for the boards of
side $5$ and $7$ are $109/32$ and $193/128$, also larger than $1/2$.

\subsubsection*{The shape of the certificate}
The shape of the certificate reflects the course of the game (\cref{tab:tree}). During the first
four rounds \Maker{} cannot force anything, and the tree is tiny: $57$ P-states up to symmetry. From
round $5$ on \Maker{} can have five stones on a line, and from round $8$ on he can have a danger pair;
the growth of the tree to its $246{,}299$ states is caused almost entirely by these forced events,
which generate $143{,}498$ and $28{,}275$ children, against $46{,}412$ children that arise from a quiet
move of \Maker. This is the difference from the $7\times7\times7$ board, where the pairing takes over
before any forced event can occur. Here a pairing cannot exist before \Breaker's seventh stone, and
every valid certificate needs at least eight rounds (\cref{prop:rmax8} below); in our certificate all
leaves occur after rounds $8$, $9$ and $10$. For $60.8\%$ of the leaves all stones of \Maker{} are revealed, so that
the leaf certificate reduces to the existence of one pairing; at the other extreme there is a single
leaf at which none of \Maker's eight stones is revealed (\cref{tab:leaf-theta}).

\begin{proposition}[A lower bound for $\Rmax$]\label{prop:rmax8}
If the root of the abstract game is \WON{} for $\Rmax$, then $\Rmax\ge 8$.
\end{proposition}

\begin{proof}
We follow the branch of the abstract game in which nothing is revealed. Let $\Pst(B,\emptyset,t)$ be a
P-state that is \WON; it is not \VOID. If it does not satisfy (P1) of \cref{def:won}, then $t<\Rmax$ and
its child \textup{(a)}, the state $\Qst(B,\emptyset,t+1)$, is \WON, since it is not \VOID. So there is
a table move $m$ for which the child \textup{(i)}, the state $\Pst(B+m,\emptyset,t+1)$, is \WON.
Starting from the root we obtain in this way \WON{} states $\Pst(B_t,\emptyset,t)$ for $t=0,1,\dots$,
until a state $\Pst(B,\emptyset,t)$ with $t\le\Rmax$ and $|B|=t$ satisfies (P1), that is,
$\LEAF(B,\emptyset,t)$ holds. We show that $t\ge 8$.

Suppose that $t\le 7$. By (L2) there is a pairing for $(B,\emptyset)$, so $t=|B|=7$ by
\cref{prop:seven}. At least $148-49$ lines are live, so \cref{cor:slack-necessary} applies and gives
$\slack(B)\ge 7$, that is, $209-2\,|\live{B}|\ge 7$: at most $101$ lines are live, and at least $47$ lines
contain a cell of $B$. But seven cells meet at most $46$ lines. If a cell of $B$ is not busy, it lies on
at most four lines (\cref{lem:degree}), and the cells of $B$ meet at most $6\cdot 7+4=46$ lines. If all
cells of $B$ are busy, three of them have the same level. Any two of these three are collinear
(\cref{lem:busy}\,(c)), and the three lines so obtained are distinct, because no line contains three
busy cells of the same level (\cref{lem:busy}\,(a),~(b)). Each of these three lines is counted at least
twice among the $7\cdot 7$ incidences of a cell of $B$ with a line, so the cells of $B$ meet at most
$49-3=46$ lines. This is a contradiction.
\end{proof}

\subsubsection*{$\LEAF$ versus $\CC$}
What \Breaker{} needs at a leaf is the statement $\CC(B,F,t)$: a pairing exists for every admissible
placement of the unrevealed stones. The certificate uses the stronger condition $\LEAF(B,F,t)$, for
three reasons. It is decided by a small number of exact augmenting-path computations
(\cref{alg:leaf}), whereas $\CC$ quantifies over all admissible sets and was decided in our
cross-checks by enumeration, by branch and bound or by an integer program, at a much higher cost. A positive
answer comes with explicit assignments from which the surplus condition can be read off
(\cref{rem:true-cert}). And $\LEAF$ is monotone in all its arguments by a short argument
(\cref{prop:leaf-monotone}), whereas $\CC$ is monotone in \Breaker's stones only for admissible $F$
(\cref{prop:C-monotone}). The price is that $\LEAF$ is not necessary for $\CC$. The two conditions
agree when all stones are revealed and when none is revealed and $t\le 10$
(\cref{rem:leaf}, \cref{prop:leaf-equals-C}), but they differ in general (\cref{ex:C-not-leaf}); in the
cross-check of \cref{ssec:cross}, $160$ of $656$ instances that satisfy $\CC$ do not satisfy $\LEAF$. A
certificate built on $\CC$ could therefore be smaller than ours. The exact enumeration confirms
$\CC(B,F,t)$ directly for the $99.55\%$ of the leaves with at most three unrevealed stones
(\cref{tab:leaves}).

\subsection{Open problems and limitations}\label{ssec:open}

\subsubsection*{The \texorpdfstring{$5\times5\times5$}{5x5x5} board}
Whether the first player wins $5\times5\times5$ tic-tac-toe remains open; it is conjectured to be a
draw~\cite{Beck2008,GolombHales2002}. We have not attempted it. The board has $109$ lines and only
$125$ cells, its center lies on $13$ lines, a threat needs only four stones, and \Breaker{} needs at
least nine stones before a pairing can exist (computed; \cref{ssec:fmt-minstones}). Cain, Raymond and
K\"allersj\"o remark that their matching approach is not feasible for this
board~\cite{CainRaymondKallersjo2025}. We do not know whether the method of the present paper,
with definitions adapted to lines of five cells, leads to a certificate of manageable size.

\subsubsection*{The number of rounds}
In our certificate \Breaker{} plays from the table for at most ten rounds. By \cref{prop:rmax8} at least
eight rounds are necessary. Whether the root is \WON{} for $\Rmax=8$ or $\Rmax=9$ is not settled: for
$\Rmax=9$ the search failed, but with truncated lists of candidate moves, so that the failure proves
nothing (\cref{ssec:observations}).

\subsubsection*{Busy cells}
The $24$ cells on the space diagonals lie on seven lines each, and all table moves of the first six
rounds are such cells. It would be natural to restrict the table to them. For $\Rmax\le 10$ this is
impossible according to an exhaustive run of the search program; the obstruction is the set of
$14$ lines in a face of the cube (\cref{ssec:observations}). This statement rests on a single program
and has not been verified independently. Whether \Breaker{} can restrict his table moves to busy cells
for some larger $\Rmax$ is open.

\subsubsection*{The size of a certificate}
The certificate was not optimized: its $93{,}164$ table entries and $106{,}723$ leaves are what the
heuristic search happened to produce. We do not know how small a valid certificate
can be, for $\Rmax=10$ or for larger $\Rmax$, nor how much would be gained by replacing $\LEAF$ by
$\CC$ or by revealing stones in a different way (\cref{rem:eager}).

\subsubsection*{Formal verification}
The proof of \cref{thm:main} has not been formally verified. Its trusted base consists of the written
proofs of \cref{sec:board,sec:pairing,sec:leaf,sec:game}, of the correctness of at least one of the
verifying programs, and of the compilers, interpreters and machines involved (\cref{ssec:trust}). Our
safeguard is redundancy: independent implementations, audits and cross-checks. A formalization of
\cref{sec:board,sec:pairing,sec:leaf,sec:game} in a proof assistant, together with a formally verified
checker for \cref{fact:H}, would remove both the written proofs and the programs from the trusted base.
The checker is small, about $700$ lines in the simplest implementations of \cref{tab:verifiers}, and
the certificate is processed in seconds by a compiled program, so this seems a realistic project.

\subsection*{Use of software tools}
The search program, the verifiers and the other programs, the computations reported in this paper and
the drafting of the manuscript were produced with the assistance of AI-based software agents, working
under the direction of the author(s), who take responsibility for the content. The certificate was
checked by independently written programs (\cref{sec:verification}).

\subsection*{Data and code availability}
The two certificate files, the source code of all programs, the outputs of the runs reported in
\cref{sec:certificate,sec:verification} and the scripts that reproduce them are contained in the
supplementary archive described in \cref{sec:repro}, which is available at [repository URL to be
added]. The SHA-256 checksums of the certificate files are given in \cref{ssec:intro-artifact}.
